%% notes-tex/033-post-query-analysis/sections/8-next.tex

\section{What the measurements change in the plan\secstatus{todo}}
\label{sec:next}

The plan's data claims survive the build. Its implementation order meets five facts of the
store that it did not anticipate, and each is a decision to take before the first run. The
recommendations below are proposals. None has been run.

\notesec{1. Choose the label of a repeated cell}
The label table's choice, the last measurement, follows from insertion order
(Section~\ref{sec:label}). The alternatives are the mean of a cell's measurements, which
is the lower-noise target, or one measurement drawn per epoch, which keeps the noise and
shows it to the model. \textit{Recommendation:} the mean, with the measurement count kept
as a column so a loss can weight by it. The change is in the label-table loop and in
nothing else, and it touches 278{,}581 cells.

\notesec{2. Decide what a vectorless cell is}
80{,}332 Hillenmeyer cells have no complete compound vector (Section~\ref{sec:environment}),
and the plan's processor raises on them. Two routes exist. The cells can be left out of
the first round, which costs 1.3 percent of the pool and removes the ten medium and
temperature conditions, which are outside the compound question, together with
tunicamycin. Or the environment context
can carry a learned vector for the absence of a compound, which keeps them and adds a
parameter whose meaning is mixed. \textit{Recommendation:} leave them out of the first
round by a read-time selection, and record the selection beside the run.

\notesec{3. Give the token the arm, or give the model the duration}
A token per source cannot carry a field that varies inside the source
(Table~\ref{tab:physical}). For Hoepfner the variation is exactly the ploidy arm, so a
token per source and arm, five tokens in all, carries it at no cost. For Hillenmeyer the
medium and the temperature stop varying once the vectorless cells of decision 2 are left
out, and the generation count remains. It is a number, so a scalar channel for the
generation count is the direct encoding. \textit{Recommendation:} both, since the first
is a vocabulary entry and the second is one masked scalar beside the dose.

\notesec{4. Read the $\gamma$ arm as ploidy and duration together}
The only data that can fit $\gamma$ changes duration with ploidy
(Section~\ref{sec:genotype}), and the two arms correlate at a fifth of their noise limit.
A learned scalar $\gamma$ fit on these cells is an estimate of the joint effect, and a
small or unstable $\gamma$ would not show that heterozygosity is unimportant.
\textit{Recommendation:} keep the fixed-identity arm as planned, and add to the learned arm
a variant in which the compound rows of the operator are conditioned on the arm, so the
round can tell a scale from an interaction. Hillenmeyer's heterozygous cells at 5 and 20
generations, 86{,}078 and 2{,}407{,}497 of them, are the one place duration varies with
ploidy held fixed, and they can be used to estimate the duration effect separately.

\notesec{5. Decide whether the loss should follow the tails}
Standardizing by the standard deviation leaves 41 to 80 percent of each source's squared
error on 2 to 3 percent of its cells (Section~\ref{sec:target}). That may be what is
wanted, since the hits are the phenotype. It is a choice either way, and it interacts with
the scale: dividing by a robust scale instead raises the weight of the heavy-tailed sources
by up to a factor of four in scale. \textit{Recommendation:} run the first round as planned
on the squared error, log the per-source loss split into cells within and beyond three
standard deviations, and hold a Huber loss as the second-round arm if the tail share of
the loss tracks the raw share.

\notesec{What needs no change}
The record population, the orientation factors, the per-source standardization, the
41-fold compound-cold evaluation with its ceilings, the choice of the feature-class
fingerprint as the default encoder, and the expectation that pooling buys a gene prior and
not more labels for the same task are all as the plan has them, now measured on the store.

\notesec{What is still missing}
The tensor layer reads only a gene name. The store carries the copy numbers, the ploidy,
the compounds, the doses and the durations in every cell's record, and the cell table
written here holds them as columns, so the processor changes of the plan's implementation
order are read-side work and need no further build.
